\worksheettitle
  {Не пропускаем разряды}
  {\modulemark{M02-R} Коррекционная карточка}

\studentnameblock

\begin{notebox}[Как выполнять]
  Сначала определите все разряды справа налево. Если единиц какого-либо
  разряда нет, запишите в этом разряде цифру \(0\).
\end{notebox}

\section{Шаг 1. Соберите число по таблице}

Соберите число \(60\,000+3\,000+40+2\).

\begin{center}
\renewcommand{\arraystretch}{1.5}
\begin{tabular}{c|ccccc}
  \toprule
  разряд & дес. тысяч & ед. тысяч & сотни & десятки & единицы \\
  \midrule
  значение & 60 000 & 3 000 & \answerblank[13mm] & 40 & 2 \\
  цифра & 6 & 3 & \answerblank[10mm] & 4 & 2 \\
  \bottomrule
\end{tabular}
\end{center}

Получилось число \answerblank[32mm]. Почему в разряде сотен нужен ноль?
\writinglines{1}

\section{Шаг 2. Заполните таблицу самостоятельно}

Представьте число \(80\,507\) по разрядам.

\begin{center}
\renewcommand{\arraystretch}{1.5}
\begin{tabular}{c|ccccc}
  \toprule
  разряд & дес. тысяч & ед. тысяч & сотни & десятки & единицы \\
  \midrule
  цифра & \answerblank[9mm] & \answerblank[9mm] & \answerblank[9mm] &
    \answerblank[9mm] & \answerblank[9mm] \\
  значение & \answerblank[14mm] & \answerblank[14mm] & \answerblank[14mm] &
    \answerblank[14mm] & \answerblank[14mm] \\
  \bottomrule
\end{tabular}
\end{center}

Разложение:
\[
  80\,507=\answerblank[88mm].
\]

\begin{practicebox}[Блиц-задание]
  \begin{enumerate}[label=\textbf{\arabic*.}]
    \item \(70\,000+600+4=\answerblank[30mm].\)
    \item \(90\,506=\answerblank[82mm].\)
  \end{enumerate}
\end{practicebox}

\section{Шаг 3. Различайте типы ошибок}

\begin{warningbox}[Ошибка в разряде цифры]
  Ученик написал:
  \[
    70\,506=70\,000+5\,000+6.
  \]
  Цифра \(5\) стоит в разряде \answerblank[25mm] и имеет значение
  \answerblank[20mm].

  Верное разложение:
  \[
    70\,506=\answerblank[82mm].
  \]
\end{warningbox}

\begin{warningbox}[Ошибка при записи]
  Ученик написал:
  \[
    80\,000+700+5=875.
  \]
  Какие разряды он пропустил? \answerblank[70mm].

  Верное число: \answerblank[32mm].
\end{warningbox}

\begin{practicebox}[Исправьте ещё две записи]
  \begin{enumerate}[label=\textbf{\arabic*.}]
    \item \(60\,090=60\,000+900.\)
    \item \(400\,000+3\,000+8=43\,008.\)
  \end{enumerate}
  \writinglines{3}
\end{practicebox}

\begin{checkpointbox}[Повторная проверка]
  \begin{enumerate}[label=\textbf{\arabic*.}]
    \item \(40\,000+500+9=\answerblank[30mm].\)
    \item \(70\,008=\answerblank[82mm].\)
    \item Объясните, почему в первом числе нужны нули в разрядах единиц тысяч
      и десятков, а во втором разложении нет слагаемых для тысяч, сотен и
      десятков.
      \writinglines{1}
  \end{enumerate}
\end{checkpointbox}
